\documentclass[11pt,a4paper]{article}
\usepackage[T1]{fontenc}
\usepackage[utf8]{inputenc}
\usepackage{lmodern}
\usepackage[margin=25mm]{geometry}
\usepackage{amsmath,amssymb,amsthm,mathtools}
\usepackage{microtype,booktabs,array,longtable}
\usepackage{xcolor,fancyhdr,enumitem}
\usepackage{hyperref,bookmark}
\definecolor{Ink}{HTML}{233240}
\definecolor{Accent}{HTML}{245A69}
\hypersetup{colorlinks=true,linkcolor=Ink,citecolor=Accent,urlcolor=Accent,
 pdftitle={Small Divisors in Surcomplex Dynamics},
 pdfauthor={Research article prepared for Vladimir Reshetnikov},
 pdfsubject={A sharp common-domain linearization trichotomy for Hahn-supported holomorphic perturbations}}
\pagestyle{fancy}\fancyhf{}
\fancyhead[L]{\small Small divisors in surcomplex dynamics}
\fancyhead[R]{\small Research article}
\fancyfoot[C]{\thepage}
\setlength{\headheight}{14pt}
\setlength{\emergencystretch}{2em}
\numberwithin{equation}{section}
\newtheorem{theorem}{Theorem}[section]
\newtheorem{lemma}[theorem]{Lemma}
\newtheorem{proposition}[theorem]{Proposition}
\newtheorem{corollary}[theorem]{Corollary}
\theoremstyle{definition}
\newtheorem{definition}[theorem]{Definition}
\newtheorem{example}[theorem]{Example}
\theoremstyle{remark}
\newtheorem{remark}[theorem]{Remark}
\newcommand{\C}{\mathbb C}
\newcommand{\R}{\mathbb R}
\newcommand{\Q}{\mathbb Q}
\newcommand{\N}{\mathbb N}
\newcommand{\Z}{\mathbb Z}
\newcommand{\No}{\mathbf{No}}
\newcommand{\K}{\mathbb K}
\newcommand{\HH}{\mathcal H}
\newcommand{\GG}{\mathcal G}
\newcommand{\RR}{\mathcal R}
\newcommand{\FF}{\mathcal F}
\newcommand{\EE}{\mathcal E}
\newcommand{\mm}{\mathfrak m}
\newcommand{\nn}{\mathfrak n}
\newcommand{\LL}{\mathcal L}
\newcommand{\id}{\operatorname{id}}
\newcommand{\supp}{\operatorname{supp}}
\newcommand{\rad}{\operatorname{rad}}
\newcommand{\st}{\operatorname{st}}
\newcommand{\ord}{\operatorname{ord}}
\newcommand{\abs}[1]{\lvert #1\rvert}
\newcommand{\norm}[1]{\lVert #1\rVert}
\newcommand{\dd}{\,\mathrm d}
\newcommand{\Gstar}{\Gamma_*}
\newcommand{\pos}{\Gamma_{>0}}
\newcommand{\doi}[1]{\href{https://doi.org/#1}{doi:#1}}
\newcommand{\arxiv}[1]{\href{https://arxiv.org/abs/#1}{arXiv:#1}}
\newcommand{\repoURL}{https://github.com/VladimirReshetnikov/Surreal}
\newcommand{\snapshot}{39f2be6667ade51bca2b45daa47e289d69c09764}

\title{\textbf{Small Divisors in Surcomplex Dynamics}\\[5pt]
\large A sharp common-domain linearization trichotomy\\
\normalsize Hahn perturbations of rotations, collapsing coefficient radii,\\
and the distinction between coefficientwise and joint analyticity}
\author{Research article prepared for Vladimir Reshetnikov}
\date{21 September 2026}

\begin{document}
\maketitle
\begin{abstract}
We investigate the Schr\"oder equation for infinitesimal Hahn-supported
perturbations of an irrational complex rotation. The analytic requirement is
that each Hahn coefficient be an ordinary holomorphic function; requiring one
common disk for all coefficients is strictly stronger than requiring a
convergent germ for each coefficient separately. For
$\lambda=e^{2\pi i\theta}$, define
$\sigma(\lambda)=\limsup_{n\to\infty}n^{-1}
\log^+\abs{\lambda^n-\lambda}^{-1}$.
We prove that $\sigma=0$ gives a unique normalized linearization on exactly
the original common disk, with an explicit support certificate and a
same-domain compositional inverse. Finite $\sigma$ always gives a
coefficient-germ linearization. Both thresholds are sharp when the class of
workspaces includes $\C((t^{\Q+\sqrt2\Q}))$.
The key obstruction is the two-monomial family
$F(z)=(\lambda+u)z+vz^2/(1-z/R)$, where $u=t$ and $v=t^{\sqrt2}$.
For $\sigma>0$, the coefficient of $u^kv$ in its unique formal linearizer has
exact convergence radius $R\exp(-(k+1)\sigma)$. Thus positive finite
$\sigma$ produces individually convergent coefficients but no common
neighborhood whatsoever. We also prove a several-variable version,
coefficientwise-entire and polynomial variants, and simultaneous
linearization for commuting perturbations. A quadratic example explains why
Hahn linearization can exist although the corresponding ordinary complex
linearizer diverges. The sharp radius-collapse result and its nonlinear
category consequences are proposed as new results; the underlying
small-divisor and Hahn-algebra tools are classical. No resolution of a named
published conjecture or independently verified priority is claimed.
\end{abstract}

\medskip
\noindent\textbf{Status.} AI-assisted research draft with complete mathematical
arguments and exact finite algebraic checks. It is not refereed or
proof-assistant verified. The novelty assessment is limited to the repository
snapshot and literature identified in Section~\ref{sec:context}.

\clearpage
\setcounter{tocdepth}{2}
\tableofcontents
\clearpage

\section{The question and the proposed contribution}\label{sec:intro}

\subsection{A dynamical problem in a genuinely different analytic category}

The classical local linearization problem asks when a holomorphic map
$f(z)=\lambda z+O(z^2)$ is conjugate to its derivative by a holomorphic change
of coordinate. The denominators $\lambda^n-\lambda$ make irrational
rotations arithmetically delicate. A surcomplex analogue must first say what
``holomorphic'' means. A Hahn series with holomorphic coefficient functions
is not the same object as an ordinary holomorphic family in a small complex
parameter, nor is it the same object as an arbitrary power series over a
non-Archimedean field.

The documents in the Surreal repository distinguish precisely these sorts
of analytic categories~\cite{Repo}. We continue that distinction rather than
identifying its inequivalent rings. Our objects have the form
\begin{equation}\label{eq:introF}
 F(z)=\lambda z+\sum_{a\in A}t^a f_a(z),\qquad
 A\subset\Gamma_{>0}\text{ well ordered},\qquad f_a(0)=0.
\end{equation}
The coefficients are ordinary functions on a disk $D_R=\{z\in\C:|z|<R\}$.
There is no ordinary convergence condition in the parameter $t$ and no
uniform bound on the family $(f_a)$. A well-ordered support and coefficientwise
finiteness, not a limiting process, govern the Hahn sums.

We seek
\begin{equation}\label{eq:schroeder}
 H\circ F=\Lambda H,\qquad
 \Lambda=F'(0),\qquad
 H(z)=z+\sum_{\gamma>0}t^\gamma h_\gamma(z),\qquad
 h_\gamma(0)=h_\gamma'(0)=0.
\end{equation}
The multiplier $\Lambda$ is allowed to have an infinitesimal part. This last
point is important: the sharp common-domain obstruction below uses a
moving multiplier, not just a perturbation whose derivative is fixed
exactly at $\lambda$.

\subsection{The arithmetic invariant}

Throughout the one-variable theory, $|\lambda|=1$ and $\lambda$ is not a root
of unity. Write $\log^+x=\max(0,\log x)$ and put
\begin{equation}\label{eq:sigma}
 \sigma(\lambda)=\limsup_{n\to\infty}
 \frac{\log^+|\lambda^n-\lambda|^{-1}}{n}\in[0,\infty].
\end{equation}
For $\lambda=e^{2\pi i\theta}$ and continued-fraction denominators $q_j$,
we prove in Section~\ref{sec:arithmetic} that
\begin{equation}\label{eq:sigmaqintro}
 \sigma(\lambda)=\limsup_j\frac{\log q_{j+1}}{q_j}.
\end{equation}
Consequently, $\sigma=0$ is the familiar subexponential small-divisor
condition, often expressed as $\log q_{j+1}=o(q_j)$. It is strictly weaker
than the Brjuno condition. The distinction between the corresponding linear
and nonlinear arithmetic conditions is classical; see, for example,
Bounemoura's discussion~\cite{Bounemoura}. The content here is what happens
when the nonlinear problem is graded by positive Hahn exponents.

\subsection{Main results, before the technical definitions}

The precise versions are Theorems~\ref{thm:lifting},
\ref{thm:collapse}, and~\ref{thm:trichotomy}.

\begin{theorem}[Overview of the trichotomy]\label{thm:overview}
Consider all maps of the form \eqref{eq:introF} over set-sized ordered
abelian value groups; counterexamples already occur for
$\Gstar=\Q+\sqrt2\Q\subset\R$.
If $\sigma(\lambda)=0$, every such map with coefficients holomorphic on
$D_R$ has a unique normalized linearizer with all coefficients holomorphic
on that same disk. If $0<\sigma(\lambda)<\infty$, every such map has a
linearizer whose coefficients are convergent germs, but some maps with
coefficients on $D_R$ have no normalized linearizer whose coefficients share
any disk $D_r$, $r>0$. If $\sigma(\lambda)=\infty$, even universal
coefficient-germ linearization fails.
In every regime, polynomial coefficient perturbations admit polynomial
coefficient linearizers.
\end{theorem}

Here ``every'' permits arbitrary well-ordered positive supports, arbitrary
coefficient sizes, and infinitesimal changes of the multiplier. It does not
mean every arbitrary germ-analytic function on $\No[i]$. Conversely, failure
of a common domain means failure within the normalized Hahn-holomorphic
category, not a claim about all imaginable set-theoretic conjugacies.

The quantitative reason for the middle regime is particularly explicit.
Set $u=t$, $v=t^{\sqrt2}$, and
\begin{equation}\label{eq:familyintro}
 F_{u,v}(z)=(\lambda+u)z+vP_R(z),\qquad
 P_R(z)=\frac{z^2}{1-z/R}.
\end{equation}
Write the unique formal solution as
$H(z)=z+\sum_{k+\ell>0}u^kv^\ell h_{k,\ell}(z)$.
For every $k\ge0$, if $\sigma>0$, we obtain
\begin{equation}\label{eq:radiusintro}
 \boxed{\quad\rad(h_{k,1})=R\exp(-(k+1)\sigma).\quad}
\end{equation}
For $0<\sigma<\infty$, every displayed radius is positive and their infimum
is zero. This is not merely loss of the input radius at the first step:
it prevents the existence of any common neighborhood after all steps.

\subsection{Why the result is not an ordinary Siegel theorem}

A Hahn coefficient at a fixed exponent depends on only finitely many
perturbation coefficients and finitely many homological inversions. There
is no demand that the resulting coefficients have a common bound as their
Hahn exponents increase. Under $\sigma=0$, each of those individual
inversions preserves the open disk. In ordinary analytic linearization one
also has to sum the entire nonlinear expansion in an ordinary variable;
that additional requirement is absent here.

This distinction gives a meaningful gain and an equally meaningful
limitation. The positive theorem gives actual maps on finite surcomplex
halos, not just a bare formal symbol, but it does not extend those maps to
arbitrary infinite spatial scales. The scale obstruction in
Section~\ref{sec:quadratic} makes the limitation visible.

\section{Analytic categories and support discipline}\label{sec:categories}

\subsection{A set-sized Hahn workspace}

Let $\Gamma$ be a set-sized ordered abelian group, written additively, and
let $\K_\Gamma=\C((t^\Gamma))$ be its Hahn field. An element is a series
$\sum_\gamma c_\gamma t^\gamma$ with well-ordered support. Its valuation is
$\nu(\sum c_\gamma t^\gamma)=\min\{\gamma:c_\gamma\ne0\}$, with
$\nu(0)=+\infty$. Put
\[
 \mathcal O_\Gamma=\{x:\nu(x)\ge0\},\qquad
 \mm_\Gamma=\{x:\nu(x)>0\}.
\]
The standard part on $\mathcal O_\Gamma$ is its coefficient at exponent
zero. It takes values in $\C$.

When $\Gamma$ is an ordered additive subgroup of $\No$, the assignment
$t^\gamma=\omega^{-\gamma}$ identifies this workspace with the corresponding
normal-form subfield of $\No[i]$. In particular, the counterexamples use only
real exponents in $\Gstar$ and have this direct surcomplex interpretation.
We need neither a set of all surcomplex numbers nor a proper-class support.
All sums, inductions, coefficient families, and value groups used in a single
construction are sets. The ordinary variable $z$ and each Taylor degree are
ordinary complex and finite-dimensional objects.

\begin{definition}[Summability]
A family $(x_i)_{i\in I}$ of Hahn series is summable if the union of its
supports is well ordered and each exponent occurs in the support of only
finitely many $x_i$. Its sum is then formed coefficient by coefficient.
\end{definition}

Summability in this definition is not convergence of partial sums in an
ordinary complex parameter. Nothing in the proofs below asserts such
convergence. We also keep two absolute-value notions apart: the ordinary
complex $|\lambda^n-\lambda|$ controls holomorphic coefficient radii, whereas
the Hahn valuation of every nonzero complex constant is zero.

\subsection{Four coefficient categories}

For a complex algebra $B$, write
\[
 B[[t^{\Gamma_{\ge0}}]]_{\mathrm H}
 =\left\{\sum_{\gamma\ge0}t^\gamma b_\gamma:
               \supp(b)\text{ is well ordered}\right\}.
\]
The subscript reminds us that the exponents need not be integers. We use
\begin{align*}
 \HH_\Gamma(R)&=\mathcal O(D_R)[[t^{\Gamma_{\ge0}}]]_{\mathrm H},\\
 \GG_\Gamma&=\C\{z\}[[t^{\Gamma_{\ge0}}]]_{\mathrm H},\\
 \FF_\Gamma&=\C[[z]][[t^{\Gamma_{\ge0}}]]_{\mathrm H},\\
 \EE_\Gamma&=\mathcal O(\C)[[t^{\Gamma_{\ge0}}]]_{\mathrm H}.
\end{align*}
Here $\C\{z\}$ is the ring of convergent ordinary power-series germs at
zero. The ring of \emph{common-domain germs} is
\[
 \RR_\Gamma=\varinjlim_{R>0}\HH_\Gamma(R),
\]
with restriction as $R$ decreases. Taylor expansion gives injections
$\RR_\Gamma\hookrightarrow\GG_\Gamma\hookrightarrow\FF_\Gamma$.
The first injection is generally not surjective: the family
$\sum_{m\ge1}t^m/(1-mz)$ has a convergent germ at every Hahn coefficient
but no common disk. This elementary distinction is emphasized in the
repository~\cite{Repo}; our obstruction will derive a much sharper dynamical
version of it.

For any of these rings, $\nn$ denotes the ideal of series supported in
$\Gamma_{>0}$. A normalized coordinate change is
$H=z+U$ with $U\in\nn$ and every coefficient of $U$ divisible by $z^2$.
Whenever the category is specified, the same requirement applies to the
coefficient functions or germs in that category.

\begin{remark}[A common disk is not a uniform estimate]
Membership in $\HH_\Gamma(R)$ requires each coefficient to be holomorphic
on $D_R$. It does not impose a common supremum bound on a smaller disk,
a norm summability condition, or a growth condition in $\gamma$. Our
positive theorem would not, as stated, prove those stronger conclusions.
\end{remark}

\subsection{The support lemma, including finite dependence}

For $A\subset\Gamma_{>0}$, let
\[
 M(A)=\{0\}\cup\{a_1+\cdots+a_k:k\ge1,\ a_i\in A\}.
\]
The following classical support principle is often called the
Hahn--Neumann lemma. We give its reduction to the finite-word theorem to
make explicit the finiteness actually used in our recursion.

\begin{lemma}[Positive supports]\label{lem:neumann}
If $A\subset\Gamma_{>0}$ is well ordered, then $M(A)$ is well ordered.
For each $\gamma\in\Gamma$, there are only finitely many finite ordered
words $(a_1,\ldots,a_k)$ in $A$ whose sum is $\gamma$. In particular,
each fixed coefficient in a Taylor expansion in a positive Hahn series
receives only finitely many contributions.
\end{lemma}

\begin{proof}
Use the finite-word well-quasi-ordering theorem of
Higman~\cite{Higman}: if an alphabet is well-quasi-ordered, its finite words
are well-quasi-ordered by subsequence embedding with componentwise
increase. A well-ordered alphabet has this property. If a word embeds in
another, its sum is at most the other's sum. Equality of sums is possible
only when no positive letter has been inserted and no letter has been
strictly increased; then the words are identical.

A strictly decreasing sequence of elements of $M(A)$, after choosing a
word for each term, would give two words in their original order with the
first embedding in the second, contradicting decreasing sums. Thus no such
sequence exists and $M(A)$ is well ordered. Infinitely many distinct words
with a fixed sum would likewise contain an embedding pair. The equality
observation makes that impossible. These arguments prove both claims.

For completeness, the finite-word theorem in the well-ordered alphabet
case has the following standard minimal-bad-sequence proof. Suppose an
infinite sequence with no embedding pair exists. Successively choose each
word of smallest length that allows such a continuation. No chosen word is
empty. Write $w_i=a_iv_i$, and select an infinite nondecreasing subsequence
of the letters $a_i$, at indices $i_1<i_2<\cdots$; every sequence in a
well-order has such a subsequence. The sequence
$w_1,\ldots,w_{i_1-1},v_{i_1},v_{i_2},\ldots$ cannot be bad, by minimality
at position $i_1$. An embedding pair among its initial words already
contradicts the original choice. An initial word embedding in a later
$v_{i_j}$ also embeds in $w_{i_j}$. Finally, an embedding of one selected
$v_{i_j}$ in a later one extends to the corresponding $w$'s because their
first letters are nondecreasing. Each possibility is a contradiction.
\end{proof}

\begin{corollary}[Finite ancestral data]\label{cor:ancestry}
Fix $\gamma\in M(A)$. Any expression formed by successively multiplying
positive-support inputs from $A$, while preserving exponent addition,
involves only finitely many input exponents and finite depths at the
coefficient $t^\gamma$.
\end{corollary}

\begin{proof}
Flatten each occurrence of an intermediate exponent into a word in $A$.
There are finitely many words of total exponent $\gamma$ by
Lemma~\ref{lem:neumann}. Their letters and lengths form finite sets.
\end{proof}

The conclusion is stronger than saying that a formal infinite sum has
well-ordered support. It rules out an infinite chain of radius losses
\emph{inside one fixed Hahn coefficient}. The radii can still decrease
without bound as the exponent itself varies; that is exactly what our
counterexample does.

\subsection{Composition and inversion without ordinary parameter limits}

If $F(z)=\lambda z+P(z)$ and $P\in\nn$, define composition by
\begin{equation}\label{eq:taylorcomposition}
 U(F(z))=
 \sum_{\beta>0}\sum_{k\ge0}
 t^\beta\frac{u_\beta^{(k)}(\lambda z)}{k!}P(z)^k,
 \qquad U=\sum_{\beta>0}t^\beta u_\beta.
\end{equation}
In the formal or germ categories this is a coefficientwise definition. In
$\HH_\Gamma(R)$ it is valid when $|\lambda|=1$, because
$\lambda D_R=D_R$. Lemma~\ref{lem:neumann}, applied to the union of the
positive supports involved, makes every Hahn coefficient a finite sum.
The same argument justifies regrouping and the usual formal Taylor
identities. These identities may also be checked at each coefficient by
passing to a finite nilpotent polynomial calculation: choose the finitely
many ancestral words and discard higher contributions.

\begin{lemma}[Near-identity inverses]\label{lem:inverse}
Let $H=z+U$ be normalized, with coefficients in any of
$\mathcal O(D_R)$, $\C\{z\}$, $\C[[z]]$, $\mathcal O(\C)$, or $\C[z]$,
and positive Hahn support $B$. There is a unique normalized inverse
$J=z+V$ with coefficient support contained in $M(B)\setminus\{0\}$ and
coefficients in that same algebra. Both $H\circ J=\id$ and
$J\circ H=\id$ hold.
\end{lemma}

\begin{proof}
In $H(J)=J+U(J)$, the coefficient $v_\gamma$ occurs with coefficient one;
every other contribution at exponent $\gamma$ uses only previously
constructed $v_\beta$, $\beta<\gamma$. Set $v_\gamma$ to minus that
finite sum and recurse on $M(B)$. Derivatives, products, and finite sums
preserve each stated coefficient algebra; the order-$z^2$ normalization
is preserved as well. This constructs a right inverse. The identical
triangular argument constructs a left inverse. Associativity, which
follows coefficientwise from finite Taylor identities, shows that these
two inverses coincide: $J_{\rm left}\circ(H\circ J_{\rm right})=
(J_{\rm left}\circ H)\circ J_{\rm right}$.
\end{proof}

\section{The ordinary homological equation}\label{sec:homological}

\subsection{An exact coefficient formula}

For $h\in z^2\C[[z]]$, put
\[
 \LL_\lambda h=h(\lambda z)-\lambda h(z),\qquad
 d_n=\lambda^n-\lambda\quad(n\ge2).
\]
Nonresonance means $d_n\ne0$. Thus $\LL_\lambda$ is a bijection on
$z^2\C[[z]]$, with
\begin{equation}\label{eq:homologicalinverse}
 \LL_\lambda^{-1}\!\left(\sum_{n\ge2}q_nz^n\right)
       =\sum_{n\ge2}\frac{q_n}{d_n}z^n.
\end{equation}
This elementary homological inverse is standard in local dynamics; its
role in nonlinear linearization is discussed in
\cite{Marmi,FauvetMenousSauzin}. We record exactly which analytic spaces
it preserves, since no other small-divisor estimate is required for the
positive Hahn theorem.

\begin{proposition}[Radius and domain preservation]\label{prop:homological}
Let $q\in z^2\C[[z]]$, let $\rho=\rad(q)$, and let $\sigma=\sigma(\lambda)$.
The following assertions hold.
\begin{enumerate}[label=(\roman*),leftmargin=2em]
\item If $\sigma=0$, the inverse $\LL_\lambda^{-1}$ maps
$z^2\mathcal O(D_R)$ into itself for every $R>0$.
\item If $\sigma<\infty$ and $0<\rho<\infty$, then
$\rad(\LL_\lambda^{-1}q)\ge\rho e^{-\sigma}$. In particular, it
preserves convergent germs.
\item If $\sigma<\infty$, it preserves entire functions.
For arbitrary $\sigma$, it preserves polynomials.
\item For $P_R(z)=z^2/(1-z/R)$,
$\rad(\LL_\lambda^{-1}P_R)=Re^{-\sigma}$,
where $e^{-\infty}=0$.
\end{enumerate}
\end{proposition}

\begin{proof}
The definition of $\sigma$ implies
\begin{equation}\label{eq:divisorroot}
 \limsup_{n\to\infty}|d_n|^{-1/n}=e^\sigma.
\end{equation}
Indeed, $|d_n|\le2$, so replacing $\log|d_n|^{-1}$ with its positive part
does not change the limsup after division by $n$.
The Cauchy--Hadamard formula applied to \eqref{eq:homologicalinverse}
gives (ii). When $\sigma=0$ and $q$ is holomorphic on $D_R$, it gives a
radius at least $R$, which proves (i) without any assertion on the
boundary. If $q$ is entire, its coefficient roots tend to zero; a finite
exponential divisor loss still leaves root limit zero. Polynomial support
in the Taylor degree is not changed at all. Finally, the coefficient of
$z^n$ in $P_R$ is $R^{2-n}$ for every $n\ge2$, so
\eqref{eq:divisorroot} gives the equality in (iv), including the
infinite-$\sigma$ case.
\end{proof}

\begin{proposition}[A useful compact-disk estimate]\label{prop:norm}
Suppose $\sigma=0$ and $0<r<s<R$. Then
\begin{equation}\label{eq:compactestimate}
 \norm{\LL_\lambda^{-1}q}_{r}
 \le \norm{q}_{s}
       \sum_{n\ge2}\frac{(r/s)^n}{|d_n|},
 \qquad
 \norm{q}_{s}=\max_{|z|\le s}|q(z)|.
\end{equation}
The numerical series on the right converges.
\end{proposition}

\begin{proof}
Cauchy's estimate gives $|q_n|\le\norm q_s s^{-n}$. Sum the resulting
absolutely convergent Taylor majorant on $|z|\le r$; its convergence follows
from \eqref{eq:divisorroot}. Thus the inverse is continuous for the
compact-open Fr\'echet topology, although an estimate on $D_r$ uses a
slightly larger compact disk. The holomorphy domain remains the full
open disk $D_R$.
\end{proof}

\subsection{Sharpness already at first order}

If $\sigma>0$, the forcing $P_R$ shows that the inverse need not preserve
$D_R$. If $\sigma=\infty$, it need not preserve convergent germs at all.
For later reference, even entire forcing can fail in the latter regime.

\begin{lemma}[Entire forcing with divergent solution]\label{lem:entirebad}
If $\sigma=\infty$, there exists $q\in z^2\mathcal O(\C)$ such that
$\LL_\lambda^{-1}q$ has radius zero.
\end{lemma}

\begin{proof}
Choose an increasing sequence $n_j$ with
$|d_{n_j}|^{-1/n_j}\to\infty$. Set
$q_{n_j}=|d_{n_j}|^{1/2}$ and all other coefficients to zero. Then
$|q_{n_j}|^{1/n_j}\to0$, so $q$ is entire, whereas
$|q_{n_j}/d_{n_j}|^{1/n_j}=|d_{n_j}|^{-1/(2n_j)}\to\infty$.
\end{proof}

These are linear facts. They do not by themselves show that a nonlinear
Hahn linearizer loses \emph{every} common neighborhood when
$0<\sigma<\infty$. Establishing that stronger failure is the purpose of
Section~\ref{sec:collapse}.

\section{Nonlinear Hahn lifting with a support certificate}\label{sec:lifting}

\subsection{The theorem in its full coefficient-algebra form}

Write
\begin{equation}\label{eq:data}
 \begin{aligned}
 F(z)&=\lambda z+P(z),&P(z)&=\sum_{a\in A}t^a f_a(z),\\
 \delta_a&=f_a'(0),&\delta&=\sum_{a\in A}t^a\delta_a,\\
 \Lambda&=\lambda+\delta,&N(z)&=P(z)-\delta z.
 \end{aligned}
\end{equation}
Thus every coefficient of $N$ vanishes to order at least two.
A coefficient indexed outside its declared support is always interpreted
as zero.

\begin{theorem}[Support-controlled linearization]\label{thm:lifting}
Assume $\lambda$ is an irrational unit multiplier and $A\subset\Gamma_{>0}$
is well ordered. For any $F$ as in \eqref{eq:data}, with
$f_a\in z\C[[z]]$, there is a unique normalized positive-support formal
solution $H=z+U$ of
\[
 H(F(z))=\Lambda H(z).
\]
Its support satisfies
\begin{equation}\label{eq:supportcertificate}
 \supp(H-z)\subset M(A)\setminus\{0\}.
\end{equation}
It has a normalized compositional inverse with the same support bound.
Moreover:
\begin{enumerate}[label=(\roman*),leftmargin=2em]
\item If $\sigma=0$ and all $f_a$ are holomorphic on $D_R$, all coefficients
of $H$ and $H^{-1}$ are holomorphic on that same $D_R$.
\item If $\sigma<\infty$ and all $f_a$ are convergent germs, the coefficients
of $H$ and $H^{-1}$ are convergent germs.
\item If $\sigma<\infty$ and all $f_a$ are entire, the coefficients of
$H$ and $H^{-1}$ are entire.
\item If all $f_a$ are polynomials, all coefficients of $H$ and $H^{-1}$
are polynomials, with no arithmetic assumption beyond nonresonance.
\end{enumerate}
No uniform coefficient bound, lower Archimedean bound on the positive
exponents, or convergence in a complex parameter is assumed.
\end{theorem}

\subsection{The recursion and its proof}

\begin{proof}
Substitute $H=z+U$ in the conjugacy equation and expand at $\lambda z$.
After cancelling $\Lambda z$, one obtains
\begin{equation}\label{eq:masterrecursion}
 \LL_\lambda U
 =-N+\delta U-
 \sum_{\beta>0}\sum_{k\ge1}
 t^\beta\frac{h_\beta^{(k)}(\lambda z)}{k!}P(z)^k.
\end{equation}
At exponent $\gamma\in M(A)\setminus\{0\}$, this becomes
\begin{equation}\label{eq:recursion}
 \begin{split}
 \LL_\lambda h_\gamma={}&-N_\gamma
 +\sum_{\substack{a\in A,\ \beta\in M(A)_{>0}\\a+\beta=\gamma}}
                 \delta_a h_\beta\\
 &-\sum_{\substack{k\ge1,\ \beta\in M(A)_{>0}\\
                  a_1,\ldots,a_k\in A\\
                  \beta+a_1+\cdots+a_k=\gamma}}
       \frac{h_\beta^{(k)}(\lambda z)}{k!}
       f_{a_1}(z)\cdots f_{a_k}(z).
 \end{split}
\end{equation}
Every sum on the right is finite. To see this even when the order on
$\Gamma$ is non-Archimedean, represent $\beta$ by a word in $A$ and append
the $a_i$'s. Lemma~\ref{lem:neumann} gives only finitely many possible
flattened words of sum $\gamma$. Their finitely many splits give only
finitely many tuples in \eqref{eq:recursion}. Every occurring $\beta$
is strictly less than $\gamma$, since all appended exponents are positive.

The right side vanishes to order at least two. The term $N_\gamma$ does so
by definition; $\delta_a h_\beta$ does so by induction. In the last line,
$h_\beta^{(k)}$ vanishes to order at least $\max(2-k,0)$, while the product
of the $k$ functions $f_{a_i}$ vanishes to order at least $k$. Their product
therefore has order at least two for $k=1$ and for every $k\ge2$.

Since $M(A)$ is well ordered and $\LL_\lambda$ is invertible on
$z^2\C[[z]]$, equation \eqref{eq:recursion} defines $h_\gamma$ by
well-founded recursion. Setting all coefficients outside $M(A)$ to zero
constructs $H$. The support is a subset of a well-ordered set, so $H$ is
an admissible Hahn section. The coefficientwise derivation of the recursion
shows that it satisfies the conjugacy equation, not merely a finite
truncation of that equation.

For (i), apply induction on $\gamma$. The operations on the right of
\eqref{eq:recursion} are differentiation, rotation, multiplication, and
finite addition on $D_R$. They preserve holomorphy on $D_R$.
Proposition~\ref{prop:homological}(i) then does the same for the
homological inverse. This induction does not successively replace the
common disk by a sequence of smaller disks. It retains exactly $D_R$
at each exponent.

For (ii), a finite sum of products of convergent germs is again a
convergent germ, and Proposition~\ref{prop:homological}(ii) applies.
There may be different resulting radii for different $\gamma$, but each
individual radius is positive. The same reasoning with entire functions
proves (iii). With polynomial coefficients, the right side is a polynomial
and the homological inverse merely divides its finitely many Taylor
coefficients by nonzero constants. This proves (iv).

It remains to verify uniqueness without assuming in advance the support
bound \eqref{eq:supportcertificate}. Let $H_1,H_2$ be two normalized
positive-support formal solutions. If they differ, the well-ordered union
of their supports contains a least exponent $\gamma$ at which their
coefficients differ. Subtract \eqref{eq:masterrecursion}. At exponent
$\gamma$, every term on the right involving a difference has a positive
extra exponent, so uses only a lower difference and is zero. Hence
$\LL_\lambda(h_{1,\gamma}-h_{2,\gamma})=0$, contradicting its
invertibility on $z^2\C[[z]]$.

Finally, Lemma~\ref{lem:inverse} constructs the inverse in the same
coefficient algebra, with support in
$M(M(A)\setminus\{0\})\subset M(A)$. This also proves
$H\circ F\circ H^{-1}(z)=\Lambda z$ as a Hahn identity.
\end{proof}

\begin{remark}[Why large Taylor coefficients are allowed]
Even if $|d_n|$ is extraordinarily small, its reciprocal is a complex
constant and has Hahn valuation zero. Polynomial coefficients therefore
survive every homological inversion regardless of their size. For
nonpolynomial coefficients, their ordinary radius is the additional
invariant that must be controlled. Neither statement can be replaced by a
claim that ``all small divisors are units, so all analytic problems disappear.''
\end{remark}

\subsection{The first two orders in a one-parameter example}

Take
\[
 F(z)=\lambda z+t f(z),\qquad
 a=f'(0),\qquad n(z)=f(z)-az,
\]
and write $H=z+t h_1+t^2h_2+\cdots$. The first equations are
\begin{align}
 \LL_\lambda h_1&=-n,\label{eq:firstorder}\\
 \LL_\lambda h_2&=a h_1-h_1'(\lambda z)f(z),\label{eq:secondorder}\\
 \LL_\lambda h_3&=a h_2-h_2'(\lambda z)f(z)
                   -\tfrac12 h_1''(\lambda z)f(z)^2.
\end{align}
These formulas display the mechanism: the new Hahn coefficient is obtained
by a finite differential expression in earlier ones and one homological
inverse. They also show why multiplier variation contributes the term
$a h_j$. Omitting that term would give the wrong recurrence and would miss
the obstruction developed below.

\subsection{Evaluation on surcomplex halos}\label{subsec:halos}

For an ordinary disk define its finite halo in the workspace by
\[
 D_R^\#=\{z_0+\zeta:z_0\in D_R,\ \zeta\in\mm_\Gamma\}.
\]
The representation $z_0+\zeta$ is unique. For
$Q=\sum_{\gamma\ge0}t^\gamma q_\gamma\in\HH_\Gamma(R)$, set
\begin{equation}\label{eq:haloeval}
 Q^\#(z_0+\zeta)=
 \sum_{\gamma\ge0}\sum_{k\ge0}
 t^\gamma\frac{q_\gamma^{(k)}(z_0)}{k!}\zeta^k.
\end{equation}
The positive supports of $Q$ and $\zeta$, together with zero, satisfy
Lemma~\ref{lem:neumann}. Thus the double family is summable. No estimate
on the sizes of the derivatives in $\gamma$ or $k$ is needed.

\begin{proposition}[Actual conjugacy on a finite halo]\label{prop:halo}
Under Theorem~\ref{thm:lifting}(i), $H^\#$ is a bijection of $D_R^\#$,
with inverse $(H^{-1})^\#$. Both $F^\#$ and multiplication by $\Lambda$
are bijections of this halo, and
\[
 H^\#\!\left(F^\#(w)\right)=\Lambda H^\#(w)
       \qquad(w\in D_R^\#).
\]
If the coefficients are entire and the hypotheses of
Theorem~\ref{thm:lifting}(iii) or (iv) apply, the same statement holds on
$\C^\#=\C+\mm_\Gamma$.
\end{proposition}

\begin{proof}
Hahn--Taylor evaluation respects composition whenever the reductions map
the base domains into one another. This follows by expanding both sides
and regrouping at a fixed exponent; Lemma~\ref{lem:neumann} reduces the
calculation to finite Taylor identities. Here the reductions of $H$ and
$H^{-1}$ are the identity, so they preserve $D_R$. The reduction of $F$
and of multiplication by $\Lambda$ is $z\mapsto\lambda z$, which preserves
$D_R$ because $|\lambda|=1$. Evaluation of the two inverse identities and
the conjugacy identity gives the result. Multiplication by $\Lambda$ has
an inverse with reduction $\lambda^{-1}$, and conjugacy then gives the
bijectivity of $F^\#$ as well.
\end{proof}

\begin{proposition}[The smaller infinitesimal monad needs no arithmetic]\label{prop:monad}
For every irrational unit multiplier, the formal coefficient solution of
Theorem~\ref{thm:lifting} can be evaluated on
$\mm_\Gamma$ and yields a bijective conjugacy there, even when some of its
ordinary coefficient series have radius zero.
\end{proposition}

\begin{proof}
At an argument $\zeta\in\mm_\Gamma$, use the formal Taylor coefficients
at zero in \eqref{eq:haloeval}. A coefficient may be an arbitrary element
of $\C[[z]]$: all its scalar coefficients still have valuation zero, and
powers of $\zeta$ have positive-support factors. The same support lemma
makes the resulting family summable. The condition $F(0)=0$ ensures that
$F(\zeta)\in\mm_\Gamma$; the normalized inverse identities and conjugacy
identity survive evaluation. This proves the assertion.
\end{proof}

Thus the negative theorems below concern analyticity of the \emph{ordinary
coefficient functions on neighborhoods}, not an absence of every form of
surcomplex linearization at infinitesimal arguments. This separation is
essential when comparing our claims with non-Archimedean dynamical results.

\section{An exact collapsing-radius obstruction}\label{sec:collapse}

\subsection{Two independent monomials in one Hahn field}

Work in the divisible ordered group
\[
 \Gstar=\Q+\sqrt2\Q\subset\R,
 \qquad u=t,\qquad v=t^{\sqrt2}.
\]
The map $(k,\ell)\mapsto k+\ell\sqrt2$ is injective on $\N^2$.
Moreover, only finitely many such pairs lie below any fixed real bound.
Consequently, the image of $\N^2$ is a well-ordered additive monoid, and
an arbitrary double formal series
$\sum_{k,\ell\ge0}b_{k,\ell}u^kv^\ell$ has an unambiguous interpretation
as a Hahn series with this support. This is a use of two independent
\emph{exponents}, not two ordinary complex parameter values.

For a fixed finite $R>0$, consider again
\[
 F_{u,v}(z)=(\lambda+u)z+vP_R(z),\qquad
 P_R(z)=\sum_{n\ge2}R^{2-n}z^n.
\]
Both perturbation coefficients are holomorphic on $D_R$ and the multiplier
is $\Lambda=\lambda+u$. Theorem~\ref{thm:lifting} gives a unique formal
linearizer with support in $\N+\sqrt2\N$.

At $v=0$ the map is already linear and its normalized linearizer is the
identity. This can be seen either by setting $v=0$ in the recursion or by
ordinary formal uniqueness for the multiplier $\lambda+u$, since every
$(\lambda+u)^n-(\lambda+u)$ has nonzero constant term in $u$.
Thus its first nontrivial $v$-block is
\begin{equation}\label{eq:Qblock}
 H(z)=z+vQ(u,z)+O(v^2),\qquad
 Q(u,z)=\sum_{k\ge0}u^k h_{k,1}(z).
\end{equation}
Here $O(v^2)$ means only the blocks with $v$-degree at least two; it is not
an asymptotic assertion in an ordinary real or complex parameter.

Extracting the coefficient of $v$ in $H(F)=\Lambda H$ gives
\begin{equation}\label{eq:Qeq}
 Q(u,(\lambda+u)z)-(\lambda+u)Q(u,z)=-P_R(z).
\end{equation}
Hence
\begin{equation}\label{eq:blockformula}
 h_{k,1}(z)=-\sum_{n\ge2}R^{2-n}c_{n,k}z^n,
 \qquad
 c_{n,k}=[u^k]\frac1{(\lambda+u)^n-(\lambda+u)}.
\end{equation}
The radius problem is reduced to an exact estimate of the reciprocal jets
$c_{n,k}$ for fixed $k$ and increasing Taylor degree $n$.

\subsection{Reciprocal jets and dominance of the highest pole}

Put
\[
 a_{n,1}=n\lambda^{n-1}-1,
 \qquad
 a_{n,j}=\binom nj\lambda^{n-j}\quad(j\ge2),
\]
where $\binom nj=0$ for $j>n$. Then
\[
 (\lambda+u)^n-(\lambda+u)
 =d_n+\sum_{j\ge1}a_{n,j}u^j.
\]
Expansion of the reciprocal in formal powers of $u$ gives, for $k\ge1$,
\begin{equation}\label{eq:reciprocaljet}
 c_{n,k}=\sum_{r=1}^k(-1)^r d_n^{-r-1}
     \sum_{\substack{j_1+\cdots+j_r=k\\j_i\ge1}}
           a_{n,j_1}\cdots a_{n,j_r},
 \qquad c_{n,0}=d_n^{-1}.
\end{equation}
For fixed $k$, the term with $r=k$ is unique and equals
\begin{equation}\label{eq:highestpole}
 (-1)^k a_{n,1}^k d_n^{-k-1}.
\end{equation}
It is this highest pole in the small divisor, rather than just its first
power, that drives the successive radius losses.

\begin{lemma}[Uniform reciprocal-jet estimates]\label{lem:jets}
For each fixed $k\ge1$ there is a constant $C_k>0$, independent of $n$,
such that
\begin{align}
 |c_{n,k}|&\le C_k n^k\max\{1,|d_n|^{-k-1}\},\label{eq:jetupper}\\
 \left|c_{n,k}-(-1)^ka_{n,1}^kd_n^{-k-1}\right|
 &\le C_k n^k|d_n|^{-k}
       \quad\text{when }|d_n|\le1.\label{eq:jetremainder}
\end{align}
In addition, $|a_{n,1}|\ge n-1$. Consequently, along any sequence with
$n\to\infty$ and $d_n\to0$,
\begin{equation}\label{eq:jetasymptotic}
 c_{n,k}=(-1)^ka_{n,1}^kd_n^{-k-1}(1+o(1)).
\end{equation}
\end{lemma}

\begin{proof}
The reverse triangle inequality gives
$|n\lambda^{n-1}-1|\ge n-1$; also $|a_{n,1}|\le n+1$.
For $j\ge2$, $|a_{n,j}|\le n^j/j!$. Therefore any product indexed by a
composition $j_1+\cdots+j_r=k$ is bounded by a constant depending only
on $k$ times $n^k$. There are finitely many compositions of $k$.
Summing \eqref{eq:reciprocaljet} gives \eqref{eq:jetupper}, because
$|d_n|\le2$ and every pole order is at most $k+1$.

After removing the $r=k$ term, all remaining pole orders are at most $k$.
For $|d_n|\le1$ this proves \eqref{eq:jetremainder}. The magnitude of
the removed term is at least $(n-1)^k|d_n|^{-k-1}$. The ratio of the
remainder bound to that lower bound is at most
$C_k(n/(n-1))^k|d_n|$, which tends to zero when $d_n\to0$.
This proves \eqref{eq:jetasymptotic} and, importantly, excludes cancellation
of the highest-pole contribution along the subsequences that control
exponential growth.
\end{proof}

\begin{theorem}[Exact coefficient-radius collapse]\label{thm:collapse}
For the family \eqref{eq:familyintro}, suppose
$\sigma(\lambda)>0$, allowing $\sigma=\infty$. Then, for every $k\ge0$,
\begin{equation}\label{eq:jetroot}
 \limsup_{n\to\infty}|c_{n,k}|^{1/n}=e^{(k+1)\sigma},
\end{equation}
and the coefficient of $u^kv$ in its normalized formal linearizer has
exact radius
\begin{equation}\label{eq:collapse}
 \rad(h_{k,1})=R e^{-(k+1)\sigma}.
\end{equation}
If $0<\sigma<\infty$, the full normalized linearizer belongs to
$\GG_{\Gstar}$ but not to $\RR_{\Gstar}$. If $\sigma=\infty$, it does
not even belong to $\GG_{\Gstar}$.
\end{theorem}

\begin{proof}
For $k=0$, \eqref{eq:jetroot} is \eqref{eq:divisorroot}.
For fixed $k\ge1$ and finite positive $\sigma$,
\eqref{eq:jetupper} gives the upper bound $e^{(k+1)\sigma}$ after taking
$n$th roots and a limsup, since $(C_kn^k)^{1/n}\to1$.
Choose a subsequence $n_j$ on which
$n_j^{-1}\log|d_{n_j}|^{-1}\to\sigma$. Because $\sigma>0$, along that
subsequence $d_{n_j}\to0$. Equation~\eqref{eq:jetasymptotic} applies, and
\[
 |c_{n_j,k}|^{1/n_j}
 =|a_{n_j,1}|^{k/n_j}|d_{n_j}|^{-(k+1)/n_j}
   |1+o(1)|^{1/n_j}\longrightarrow e^{(k+1)\sigma}.
\]
We used $n_j-1\le|a_{n_j,1}|\le n_j+1$ for the first factor. This proves
the matching lower bound.

If $\sigma=\infty$, choose $n_j$ so that
$|d_{n_j}|^{-1/n_j}\to\infty$. Again $d_{n_j}\to0$ and the same
asymptotic implies $|c_{n_j,k}|^{1/n_j}\to\infty$. Thus
\eqref{eq:jetroot} remains true with an infinite right side.
The Cauchy--Hadamard formula in \eqref{eq:blockformula} now gives
\eqref{eq:collapse}.

When $0<\sigma<\infty$, Theorem~\ref{thm:lifting}(ii) says that every
coefficient of the full nonlinear linearizer is a convergent germ.
However, the particular radii in \eqref{eq:collapse} have infimum zero,
so its coefficients cannot all be holomorphic on any $D_r$, $r>0$.
When $\sigma=\infty$, even $h_{0,1}$ has radius zero and is not a
convergent germ. Formal uniqueness in Theorem~\ref{thm:lifting} shows
that no alternative normalized solution, even with additional positive
Hahn exponents, can avoid these coefficients.
\end{proof}

\subsection{What the example proves, and why two exponents matter}

The first coefficient $h_{0,1}$ alone only loses the radius by $e^{-\sigma}$.
For finite positive $\sigma$, that loss is compatible with the existence
of some smaller common disk. The higher multiplier jets $h_{k,1}$ force
losses by every power $e^{-(k+1)\sigma}$ and exclude such a disk. The
argument is therefore strictly stronger than a first-order homological
counterexample.

The independence of $1$ and $\sqrt2$ prevents different pairs $(k,\ell)$
from merging into the same Hahn coefficient. It lets us isolate the
$v$-linear block without contamination by higher $v$-powers. Replacing
$v$ by an integer power of $u$ would merge blocks and could introduce
cancellation; the proof above would no longer establish the claimed
radius for a particular Hahn coefficient. We do not silently make that
replacement.

\begin{corollary}[A finite-support instability]\label{cor:finitefailure}
For every irrational multiplier with $0<\sigma<\infty$, two nonzero
positive Hahn coefficients suffice for a holomorphic perturbation of a
rotation to have a coefficient-germ linearizer but no common-domain
linearizer. The input coefficients can be rational functions of $z$.
\end{corollary}

\begin{proof}
The only nonzero perturbation coefficients are $z$ at exponent $1$ and
$P_R(z)$ at exponent $\sqrt2$. Apply Theorem~\ref{thm:collapse}.
\end{proof}

\section{Sharp category thresholds}\label{sec:thresholds}

\subsection{A precise universal formulation}

Let $\mathsf{CD}(\lambda)$ be the assertion that, for every set-sized
ordered abelian group $\Gamma$, every $R>0$, and every normalized input
$F$ as in \eqref{eq:data} in $\HH_\Gamma(R)$, there is a normalized
linearizer in $\RR_\Gamma$. Let $\mathsf{CG}(\lambda)$ be the assertion
that every such input with coefficient germs has a normalized linearizer
in $\GG_\Gamma$. These are universal statements about specified
coefficient categories.

\begin{theorem}[The common-domain/coefficient-germ trichotomy]\label{thm:trichotomy}
For an irrational unit multiplier $\lambda$,
\begin{equation}\label{eq:equiv}
 \mathsf{CD}(\lambda)\ \Longleftrightarrow\ \sigma(\lambda)=0,
 \qquad
 \mathsf{CG}(\lambda)\ \Longleftrightarrow\ \sigma(\lambda)<\infty.
\end{equation}
The same equivalences hold when the universal quantifier over $\Gamma$
is replaced by the single group $\Gstar=\Q+\sqrt2\Q$. In the positive
common-domain case, the linearizer and its inverse preserve the original
disk, not just an unspecified smaller disk.
\end{theorem}

\begin{proof}
The positive implications are Theorem~\ref{thm:lifting}(i) and (ii).
For $0<\sigma<\infty$, Theorem~\ref{thm:collapse} disproves
$\mathsf{CD}(\lambda)$ within $\Gstar$ while the positive coefficient-germ
implication still holds. For $\sigma=\infty$, the same example has a
nonconvergent coefficient and disproves both universal properties.
\end{proof}

\begin{corollary}[The entire-coefficient threshold]\label{cor:entirethreshold}
Universal linearization of entire-coefficient perturbations by
entire-coefficient normalized coordinate changes holds exactly when
$\sigma<\infty$. In contrast, universal polynomial-coefficient
linearization holds for every irrational unit multiplier.
\end{corollary}

\begin{proof}
The positive assertions follow from Theorem~\ref{thm:lifting}(iii)--(iv).
For $\sigma=\infty$, choose the entire forcing $q$ from
Lemma~\ref{lem:entirebad} and set $F(z)=\lambda z+tq(z)$.
The coefficient of $t$ in its normalized linearizer is
$-\LL_\lambda^{-1}q$, which is not even a convergent germ.
\end{proof}

\subsection{The meaning of sharpness}

Table~\ref{tab:trichotomy} summarizes the result. The polynomial column
means that each Hahn coefficient is a polynomial; it does not mean that
the complete coordinate change is a finite polynomial over the Hahn field.

\begin{table}[ht]
\centering\small
\begin{tabular}{@{}>{\raggedright\arraybackslash}p{.18\textwidth}>{\raggedright\arraybackslash}p{.26\textwidth}>{\raggedright\arraybackslash}p{.24\textwidth}>{\raggedright\arraybackslash}p{.22\textwidth}@{}}
\toprule
Arithmetic & Common ordinary disk & Coefficient germs / entire coefficients
& Polynomial coefficients\\
\midrule
$\sigma=0$ & Always, with the same $D_R$ & Always / always & Always\\[3pt]
$0<\sigma<\infty$ & Not universally; explicit radii collapse to zero
& Always / always & Always\\[3pt]
$\sigma=\infty$ & Not universally & Neither universally & Always\\
\bottomrule
\end{tabular}
\caption{Universal normalized linearization in the specified Hahn categories.
Negative common-domain assertions use $\Gstar$; positive assertions hold
for every set-sized ordered abelian group.}\label{tab:trichotomy}
\end{table}

The distinction between the middle two columns is not contradictory.
For entire input coefficients, a finite exponential loss of Taylor
coefficient growth leaves every coefficient entire. For coefficients
with a finite input radius, repeated losses can exhaust every positive
common radius. The fixed coefficient algebra, not just the existence of
Hahn sums, determines the answer.

\section{Arithmetic examples and the Brjuno comparison}\label{sec:arithmetic}

\subsection{Continued fractions measure the exponential loss}

Let $\theta\in\R\setminus\Q$, let $p_j/q_j$ be its continued-fraction
convergents, and let $\|x\|_{\Z}$ denote distance to the nearest integer.
We use the classical approximation estimates
\begin{equation}\label{eq:cfbounds}
 \frac1{q_j+q_{j+1}}<\|q_j\theta\|_{\Z}<\frac1{q_{j+1}},
 \qquad
 \|m\theta\|_{\Z}\ge\|q_j\theta\|_{\Z}
       \quad(q_j\le m<q_{j+1}).
\end{equation}
These standard continued-fraction facts and their role in small-divisor
problems are treated in Marmi's lectures~\cite{Marmi}. The short argument
below translates them to our particular invariant; it is included to make
all arithmetic examples explicit.

\begin{proposition}[Continued-fraction formula]\label{prop:continued}
For $\lambda=e^{2\pi i\theta}$,
\[
 \sigma(\lambda)=\limsup_{j\to\infty}\frac{\log q_{j+1}}{q_j}.
\]
\end{proposition}

\begin{proof}
For every real $x$,
\[
 4\|x\|_{\Z}\le|1-e^{2\pi ix}|\le2\pi\|x\|_{\Z}.
\]
The upper bound follows from $|\sin y|\le|y|$, and the lower bound
from concavity of $\sin$ on $[0,\pi/2]$. Constants do not affect a
logarithmic limsup divided by $m$. Thus $\sigma$ is the same as
\[
 \limsup_{m\to\infty}\frac{\log\|m\theta\|_{\Z}^{-1}}m,
\]
since $m=n-1$ and the ratio $m/(m+1)$ tends to one. At $m=q_j$,
\eqref{eq:cfbounds} gives the lower bound
$\log q_{j+1}/q_j$ for this limsup. If $q_j\le m<q_{j+1}$, the other
half of \eqref{eq:cfbounds} gives
\[
 \frac{\log\|m\theta\|_{\Z}^{-1}}m
 \le\frac{\log(q_j+q_{j+1})}{q_j}
 \le\frac{\log q_{j+1}}{q_j}+\frac{\log2}{q_j}.
\]
This proves the reverse bound, including the possibility of an infinite
limsup.
\end{proof}

\subsection{A strictly weaker condition than Brjuno}

The Brjuno condition is
\begin{equation}\label{eq:brjuno}
 B(\theta)=\sum_{j\ge0}\frac{\log q_{j+1}}{q_j}<\infty.
\end{equation}
It implies $\sigma=0$, since the nonnegative summands tend to zero.
It is a much stronger requirement than mere convergence of those summands
to zero. The following construction gives a concrete infinite
continued-fraction specification witnessing the difference.

\begin{example}[Subexponential, but not Brjuno]\label{ex:nonbrjuno}
Start with a positive initial continued-fraction coefficient, for example
$a_1=2$, and choose successively
\begin{equation}\label{eq:cfzero}
 a_{j+1}=\left\lceil\exp\left(\frac{q_j}{j+1}\right)\right\rceil
       \qquad(j\ge1).
\end{equation}
Let $\theta=[0;a_1,a_2,\ldots]$. Then $\sigma(e^{2\pi i\theta})=0$
but $B(\theta)=\infty$.
\end{example}

\begin{proof}
The denominator recurrence is $q_{j+1}=a_{j+1}q_j+q_{j-1}$.
With $x_j=q_j/(j+1)$, the bounds
$e^{x_j}\le a_{j+1}\le e^{x_j}+1$ and $q_{j-1}\le q_j$ imply
\[
 \frac1{j+1}+\frac{\log q_j}{q_j}
 \le\frac{\log q_{j+1}}{q_j}
 \le\frac1{j+1}+\frac{\log q_j}{q_j}+\frac{\log3}{q_j}.
\]
The denominators tend to infinity, so both outside expressions tend to
zero. Proposition~\ref{prop:continued} gives $\sigma=0$. The lower bound
$1/(j+1)$ makes the Brjuno sum diverge.
\end{proof}

By Theorem~\ref{thm:lifting}, this multiplier linearizes every positive
Hahn perturbation with common-disk coefficients in the common-disk sense.
In contrast, the classical Brjuno--Yoccoz theorem implies that its ordinary
quadratic germ $w\mapsto\lambda w+w^2$ is not holomorphically
linearizable~\cite{Yoccoz}. There is no contradiction: the positive Hahn
perturbation $\lambda z+tz^2$ and the ordinary quadratic germ have different
coefficient categories and different spatial scales.

\subsection{Every remaining regime occurs}

\begin{example}[A prescribed finite loss exponent]\label{ex:positive}
For any real $s>0$, choose
$a_{j+1}=\lceil e^{s q_j}\rceil$. Then the resulting irrational
$\theta$ satisfies $\sigma(e^{2\pi i\theta})=s$.
Indeed,
\[
 s+\frac{\log q_j}{q_j}
 \le\frac{\log q_{j+1}}{q_j}
 \le s+\frac{\log q_j}{q_j}+\frac{\log3}{q_j},
\]
so the quotient tends to $s$.
\end{example}

\begin{example}[Superexponential small divisors]\label{ex:infinite}
Choose $a_{j+1}=\lceil e^{q_j^2}\rceil$. The same argument gives
\[
 \frac{\log q_{j+1}}{q_j}=q_j+o(1),
\]
so $\sigma=\infty$.
\end{example}

At the other end, any polynomial lower bound
$|\lambda^n-\lambda|\ge Cn^{-\tau}$ implies $\sigma=0$ immediately.
The preceding continued-fraction constructions are useful because they
show that all three regimes of Theorem~\ref{thm:trichotomy} occur, and that
the positive regime genuinely extends beyond the ordinary Brjuno class.

\section{Polynomial perturbations and the critical spatial scale}\label{sec:quadratic}

\subsection{An explicit coefficientwise-entire conjugacy}

For the ordinary formal quadratic map
$p_\lambda(w)=\lambda w+w^2$, let
\[
 h_\lambda(w)=w+\sum_{n\ge2}b_n w^n
\]
be its unique normalized formal solution of
$h_\lambda(p_\lambda(w))=\lambda h_\lambda(w)$.
It exists for every nonresonant $\lambda$, by induction in the ordinary
Taylor degree. More explicitly,
\begin{equation}\label{eq:quadraticrecurrence}
 b_1=1,\qquad
 b_n=-\frac{1}{\lambda^n-\lambda}
       \sum_{m=\lceil n/2\rceil}^{n-1}
         b_m\binom m{n-m}\lambda^{2m-n}
       \quad(n\ge2).
\end{equation}
The binomial coefficient is the coefficient of $w^n$ in
$(\lambda w+w^2)^m$. Thus this recurrence follows directly by comparing
Taylor coefficients and is not a convergence assertion.
The first values are
\[
 b_2=-\frac1{\lambda^2-\lambda},\qquad
 b_3=\frac{2\lambda}{(\lambda^2-\lambda)(\lambda^3-\lambda)}.
\]

Now take the positive Hahn perturbation
\begin{equation}\label{eq:quadraticF}
 F_t(z)=\lambda z+tz^2.
\end{equation}
Its normalized linearizer is
\begin{equation}\label{eq:quadraticH}
 H_t(z)=z+\sum_{n\ge2}b_n t^{n-1}z^n.
\end{equation}
At every Hahn exponent there is only one polynomial coefficient. Therefore
$H_t$ lies in the entire-coefficient category for every irrational
multiplier, including those with $\sigma=\infty$.

\begin{proposition}[Quadratic scaling identity]\label{prop:quadraticscale}
As formal identities,
\[
 H_t(z)=t^{-1}h_\lambda(tz),\qquad
 H_t(F_t(z))=\lambda H_t(z).
\]
The change of variable $H_t$ and its inverse are well defined by Hahn sums
on the full set
\begin{equation}\label{eq:subcritical}
 \{z\in\K_\Gamma:\nu(tz)>0\},
\end{equation}
whenever the workspace contains the monomial $t$ of positive valuation.
They give a conjugacy on that set.
\end{proposition}

\begin{proof}
The first identity is coefficientwise algebra. The relation
$tF_t(z)=p_\lambda(tz)$ then proves the second by formal substitution.
At any argument with $tz\in\mm_\Gamma$, the ordinary formal series
$h_\lambda(tz)$ is Hahn summable: its scalar coefficients have valuation
zero, and the powers of the positive-support argument satisfy
Lemma~\ref{lem:neumann}. Its ordinary formal inverse has the same property.
Both preserve the infinitesimal monad, because they have linear term the
identity and higher terms of strictly higher valuation. After multiplication
by $t^{-1}$, they preserve \eqref{eq:subcritical}. The identities and their
inverse counterparts may be evaluated there coefficientwise.
\end{proof}

For $t=\omega^{-1}$, the domain includes every finite surcomplex number
and many infinite ones, but it excludes the critical scale where $tz$ has
nonzero ordinary standard part. This is consistent with the familiar
non-Archimedean linearization phenomenon over a trivially valued
coefficient field; it is not presented as a new replacement for the
ultrametric results of Lindahl~\cite{Lindahl}.

\subsection{Why a divergent classical germ does not become convergent}

Rescale the spatial variable to the critical scale, $z=t^{-1}w$.
As a formal Taylor series in $w$ over the Hahn field,
\begin{equation}\label{eq:critical}
 tH_t(t^{-1}w)=h_\lambda(w).
\end{equation}
On the right, every coefficient has Hahn exponent zero. If
$\rad(h_\lambda)=0$, it is not a convergent ordinary germ.

\begin{proposition}[A precise rescaled-germ obstruction]\label{prop:critical}
The formal germ in \eqref{eq:critical} is the Taylor germ of a
common-domain Hahn-holomorphic section at $w=0$ if and only if the ordinary
formal series $h_\lambda$ has positive convergence radius.
\end{proposition}

\begin{proof}
Suppose a section $A(w)=\sum_{\gamma\ge0}t^\gamma a_\gamma(w)$, with all
$a_\gamma$ holomorphic on one $D_r$, has Taylor series equal to
$h_\lambda(w)$ in $\K_\Gamma[[w]]$. The coefficient at Hahn exponent
zero has ordinary Taylor series $h_\lambda$. It must therefore converge
at least on $D_r$. All other coefficient germs have zero Taylor series.
Conversely, if $h_\lambda$ converges on a disk, the section with only
exponent-zero coefficient $h_\lambda$ is the desired section.
\end{proof}

For the non-Brjuno multiplier of Example~\ref{ex:nonbrjuno}, Yoccoz's
quadratic theorem gives $\rad(h_\lambda)=0$~\cite{Yoccoz}.
Thus $H_t$ is coefficientwise entire and an actual conjugacy on the
subcritical domain, while its rescaled germ cannot be realized by a
common-domain analytic section at the critical scale.

Even when $h_\lambda$ converges ordinarily, substituting a nonzero ordinary
$w$ directly into the displayed Hahn family in \eqref{eq:critical} is
\emph{not} a Hahn summation: infinitely many nonzero terms would occupy
exponent zero. One must first use the separately justified ordinary
holomorphic germ, and then its canonical Taylor lift. This avoids an
illegitimate exchange of ordinary analytic summation with Hahn summation.

\subsection{What this example does not say}

A coefficientwise-entire section need not be a globally analytic map on
all of $\No[i]$ under every spatial rescaling. The common-domain category
controls ordinary coefficient domains around ordinary centers. It does
not itself supply coherence between every possible infinite-scale chart.
Likewise, the existence of the Hahn conjugacy does not imply ordinary
holomorphic dependence on $t$ near $t=0$. For a divergent $h_\lambda$, the
formal dependence encoded in \eqref{eq:quadraticH} need not converge as an
ordinary two-variable series in $(t,z)$.

\section{Commuting perturbations and simultaneous linearization}\label{sec:centralizer}

The existence theorem also gives a clean dynamical consequence. This
corollary uses the standard nonresonant centralizer argument; its role is
to show that the support and common-domain conclusions are preserved when
one passes from a single map to a commuting family.

\begin{theorem}[A common coordinate for the centralizer]\label{thm:centralizer}
Suppose $\sigma(\lambda)=0$ and $F$ satisfies
Theorem~\ref{thm:lifting}(i) on $D_R$. Let $H$ be its normalized
linearizer and $\Lambda=F'(0)$. Let
\[
 G(z)=\mu z+\sum_{b\in B}t^b g_b(z),\qquad
 |\mu|=1,\quad g_b\in z\mathcal O(D_R),\quad B\subset\Gamma_{>0}
\]
be another positive-support perturbation, and suppose $G\circ F=F\circ G$.
Then, with $M=G'(0)$,
\begin{equation}\label{eq:simultaneous}
 H\circ G\circ H^{-1}(z)=Mz.
\end{equation}
No nonresonance or small-divisor assumption on $\mu$ is required.
Consequently, every such map commuting with $F$ is linearized by the same
$H$, and any two such maps commute with each other.
\end{theorem}

\begin{proof}
All compositions are legitimate on $D_R$ by the support lemma and
$|\lambda|=|\mu|=1$. Put $K=H\circ G\circ H^{-1}$. It commutes with
$z\mapsto\Lambda z$, fixes zero, and has derivative $M$ at zero.
Expand its ordinary Taylor germ over the Hahn field:
$K(z)=\sum_{n\ge1}k_nz^n$. Comparing coefficients in
$K(\Lambda z)=\Lambda K(z)$ gives
\[
 (\Lambda^n-\Lambda)k_n=0\quad(n\ge2).
\]
The factor $\Lambda^n-\Lambda$ is a nonzero Hahn-field element: its
exponent-zero coefficient is $\lambda^n-\lambda\ne0$. Hence every
$k_n$ with $n\ge2$ vanishes and $k_1=M$. At each Hahn exponent the
corresponding holomorphic coefficient function has the Taylor series of
a linear function, so the identity theorem on $D_R$ yields $K(z)=Mz$
as a common-domain section. Equation~\eqref{eq:simultaneous} follows.
Two scalar multiplications commute; conjugating back proves the final
assertion.
\end{proof}

\begin{corollary}[No nontrivial tangent-to-identity centralizer element]
Under the preceding hypotheses, a commuting perturbation $G$ with
$G'(0)=1$ is the identity. More generally, within this class the multiplier
$G'(0)$ determines $G$ uniquely.
\end{corollary}

\begin{proof}
In the coordinate $H$, the map is multiplication by its multiplier.
\end{proof}

The same proof works in the coefficient-germ category when $\sigma<\infty$,
and in the polynomial-coefficient setting whenever the constructed
linearizers have polynomial coefficients. In those statements equality is
interpreted in the indicated category, rather than on a disk that has not
been shown to exist. A set-indexed family can be treated one member at a
time using the same $H$; no unsupported sum over the whole family is taken.

\section{Several variables: the same sharp mechanism}\label{sec:several}

\subsection{Diagonal unitary reduction and nonresonance}

Let $d\ge1$, let
\[
 L=\operatorname{diag}(\lambda_1,\ldots,\lambda_d),\qquad
 |\lambda_j|=1,
\]
and assume the nonresonance condition
\begin{equation}\label{eq:multinonresonance}
 \lambda^\alpha\ne\lambda_j
 \quad\text{for every }\alpha\in\N^d,\ |\alpha|\ge2,
 \quad j=1,\ldots,d,
 \qquad\lambda^\alpha=\lambda_1^{\alpha_1}\cdots\lambda_d^{\alpha_d}.
\end{equation}
For $n\ge2$ put
\begin{equation}\label{eq:multisigma}
 \Delta_n=\min_{\substack{|\alpha|=n\\1\le j\le d}}
                 |\lambda^\alpha-\lambda_j|>0,
 \qquad
 \sigma_d(L)=\limsup_{n\to\infty}\frac{\log^+\Delta_n^{-1}}n.
\end{equation}
The common domain is the equal-radius polydisk
$\mathbb D_R^d=\{z:\max_j|z_j|<R\}$. Our assertions are about this
specified geometry; no claim about an optimal general Reinhardt domain
is needed.

For a vector formal power series $q(z)=\sum_{|\alpha|\ge2}q_\alpha z^\alpha$,
define its isotropic radius by
\begin{equation}\label{eq:isotropic}
 \rho(q)^{-1}=
 \limsup_{n\to\infty}
 \left(\max_{\substack{|\alpha|=n\\1\le j\le d}}
              |q_{\alpha,j}|\right)^{1/n},
\end{equation}
with the usual zero/infinity conventions. It is the supremum of radii
of equal-radius polydisks of Taylor convergence. Indeed, the number of
multiindices of degree $n$ is $\binom{n+d-1}{d-1}$, whose $n$th root
tends to one. Coefficient bounds and the root test therefore establish
both implications in this interpretation.

\begin{lemma}[Multivariable homological inverse]\label{lem:multihomological}
The operator $\LL_Lh=h(Lz)-Lh(z)$ is invertible on vector formal series
of order at least two, with
\begin{equation}\label{eq:multihomological}
 (\LL_L^{-1}q)_{\alpha,j}
       =\frac{q_{\alpha,j}}{\lambda^\alpha-\lambda_j}.
\end{equation}
If $\sigma_d=0$, it preserves holomorphy on every
$\mathbb D_R^d$. If $\sigma_d<\infty$, it preserves convergent germs and
entire functions; a finite nonzero isotropic radius loses at most the
factor $e^{-\sigma_d}$. It preserves polynomial vectors under
nonresonance alone.
\end{lemma}

\begin{proof}
The monomial $z^\alpha e_j$ is an eigenvector of $\LL_L$ with eigenvalue
$\lambda^\alpha-\lambda_j$. This proves the formal formula.
Taking maxima at total degree $n$ in \eqref{eq:multihomological} yields
\[
 \max_{|\alpha|=n,j}|(\LL_L^{-1}q)_{\alpha,j}|
 \le\Delta_n^{-1}\max_{|\alpha|=n,j}|q_{\alpha,j}|.
\]
Now use \eqref{eq:isotropic}, or Cauchy estimates on nested polydisks.
The number of degree-$n$ monomials has only polynomial growth, so it adds
no exponential loss. The entire and polynomial assertions follow exactly
as in Proposition~\ref{prop:homological}.
\end{proof}

\subsection{A nonlinear common-polydisk theorem}

\begin{theorem}[Multivariable support-controlled lifting]\label{thm:multilifting}
Let
\[
 F(z)=Lz+\sum_{a\in A}t^a f_a(z),\qquad f_a(0)=0,
 \qquad M=DF(0)=L+\sum_{a\in A}t^a Df_a(0),
\]
where $A\subset\Gamma_{>0}$ is well ordered. Under
\eqref{eq:multinonresonance}, there is a unique normalized formal
coordinate change $H=z+U$, with $DU(0)=0$ coefficientwise, such that
\begin{equation}\label{eq:multischroeder}
 H(F(z))=M H(z),\qquad
 \supp(H-z)\subset M(A)\setminus\{0\}.
\end{equation}
If $\sigma_d=0$ and all $f_a$ are holomorphic on $\mathbb D_R^d$, then
$H$ and its inverse have coefficients holomorphic on that same polydisk.
For finite $\sigma_d$, coefficient germs and entire coefficients are
preserved. Polynomial coefficients are preserved for every nonresonant
$L$. The infinitesimal matrix $M-L$ need not be diagonal.
\end{theorem}

\begin{proof}
Write $P=F-Lz$, $E=M-L$, and $N=P-Ez$. The vector version of
\eqref{eq:masterrecursion} is
\begin{equation}\label{eq:multimaster}
 \LL_LU=-N+EU-
 \sum_{\beta>0}\sum_{k\ge1}\frac{t^\beta}{k!}
       D^kh_\beta(Lz)[P(z),\ldots,P(z)],
\end{equation}
where $D^kh$ is the symmetric $k$-linear derivative. At a fixed Hahn
exponent only finitely many terms occur by
Lemma~\ref{lem:neumann}, and every previously occurring $h_\beta$ has
strictly smaller exponent. The right side has total Taylor order at
least two: $P$ has order one, $h_\beta$ has order two, and the same
order count as in the one-variable proof applies to its $k$th derivative.
Apply \eqref{eq:multihomological} by well-founded recursion.

The operations in \eqref{eq:multimaster} preserve the common polydisk
because $L$ preserves it, and Lemma~\ref{lem:multihomological} provides
the required analytic inverse. For germs, entire functions, and
polynomials, use the corresponding parts of that lemma. The least-exponent
argument proves uniqueness. The inverse-coordinate recursion is vectorial
but still has coefficient one on the new unknown, so the proof of
Lemma~\ref{lem:inverse} applies with multilinear derivatives and gives
the same coefficient algebra and support certificate.
\end{proof}

As in Proposition~\ref{prop:halo}, common-polydisk solutions evaluate to
actual conjugacies on
$(\mathbb D_R^d)^\#=\mathbb D_R^d+\mm_\Gamma^d$. Both the reduction of
$F$ and the reduction of its linear normal form are $L$, which preserves
the base polydisk. Positive perturbations of the invertible matrix $L$
have a Hahn inverse, so the conjugate maps are bijections of the halo.

\subsection{A sharp multivariable radius-collapse family}

Set
\begin{equation}\label{eq:multiP}
 P_R^{(d)}(z)=
 \prod_{j=1}^d(1-z_j/R)^{-1}-1-\sum_{j=1}^d z_j/R
 =\sum_{|\alpha|\ge2}R^{-|\alpha|}z^\alpha,
\end{equation}
and let $\mathbf1=(1,\ldots,1)^T$. In the same two-monomial workspace,
consider
\begin{equation}\label{eq:multiF}
 F_{u,v}(z)=(1+u)Lz+vP_R^{(d)}(z)\mathbf1.
\end{equation}
The choice of a \emph{multiplicative} detuning $(1+u)L$ makes the highest
reciprocal pole uniform over every component and every multiindex.

\begin{theorem}[Exact multivariable radius collapse]\label{thm:multicollapse}
Assume \eqref{eq:multinonresonance} and $\sigma_d(L)>0$. In the unique
normalized formal linearizer of \eqref{eq:multiF}, let $h_{k,1}(z)$ be
the coefficient vector of $u^kv$. Then, for every $k\ge0$,
\begin{equation}\label{eq:multicollapse}
 \rho(h_{k,1})=R e^{-(k+1)\sigma_d(L)}.
\end{equation}
Consequently, the universal common-polydisk-germ property holds exactly
when $\sigma_d=0$, and the universal coefficient-germ property holds
exactly when $\sigma_d<\infty$, when workspaces include $\Gstar$.
\end{theorem}

\begin{proof}
At $v=0$, the map is linear and its normalized coordinate change is the
identity. The coefficient of $v$ therefore satisfies the homological
equation for the diagonal multiplier $(1+u)L$. In component $j$ and
at multiindex $\alpha$, with $n=|\alpha|$, its coefficient is
\[
 -\frac{R^{-n}}{(1+u)^n\lambda^\alpha-(1+u)\lambda_j}.
\]
Write
\[
 D_{\alpha,j}=\lambda^\alpha-\lambda_j,
 \qquad
 A_{\alpha,j,1}=n\lambda^\alpha-\lambda_j,
 \qquad
 A_{\alpha,j,r}=\binom nr\lambda^\alpha\quad(r\ge2).
\]
The reciprocal-jet formula \eqref{eq:reciprocaljet} remains valid with
these replacements. Because all the phases have modulus one,
\[
 |A_{\alpha,j,1}|\ge n-1,
 \qquad |A_{\alpha,j,r}|\le n^r/r!\quad(r\ge2).
\]
Thus, for each fixed $k$, the estimates in Lemma~\ref{lem:jets} hold
with constants independent not only of $n$ but also of $\alpha$ and $j$.

Let $C_{\alpha,j,k}$ be the corresponding coefficient of $u^k$ in the
reciprocal. The upper estimate yields
\[
 \max_{|\alpha|=n,j}|C_{\alpha,j,k}|
 \le C_k n^k\max\{1,\Delta_n^{-k-1}\}.
\]
For the lower bound, choose $n$ along a subsequence realizing the
limsup in \eqref{eq:multisigma}, and at each such degree choose
$(\alpha,j)$ attaining the finite minimum $\Delta_n$. Positive
$\sigma_d$ ensures $\Delta_n\to0$ along this subsequence.
The highest-pole term dominates its remainder, uniformly in the chosen
multiindices, and therefore
\[
 \limsup_{n\to\infty}
 \left(\max_{|\alpha|=n,j}|C_{\alpha,j,k}|\right)^{1/n}
       =e^{(k+1)\sigma_d}.
\]
For infinite $\sigma_d$ the same subsequence argument gives an infinite
limsup. Multiplication by $R^{-n}$ and the definition
\eqref{eq:isotropic} prove \eqref{eq:multicollapse}.

For finite positive $\sigma_d$, Theorem~\ref{thm:multilifting} ensures
that every coefficient of the full linearizer is a convergent germ.
But the displayed isotropic radii have infimum zero, so no common
ordinary neighborhood exists: every neighborhood of zero in $\C^d$
contains an equal-radius polydisk. For infinite $\sigma_d$, the first
$v$ coefficient already fails to be a germ. The positive threshold
implications were proved in Theorem~\ref{thm:multilifting}, and uniqueness
excludes alternative positive-support solutions.
\end{proof}

\begin{remark}[Nonresonance is indispensable]
If $\lambda^\alpha=\lambda_j$, a first-order forcing
$t z^\alpha e_j$ lies outside the range of $\LL_L$ at that monomial.
It cannot be removed by a normalized near-identity coordinate change.
Resonant normal forms require a different target equation, not the
linear normal form used in this article. The multivariable result does
not presume that every tuple of unit multipliers is nonresonant.
\end{remark}

\section{Relation to the repository and published work}\label{sec:context}

\subsection{The repository snapshot and the gap selected}

The initial repository inspection was pinned to commit
\begin{center}
\small\texttt{39f2be6667ade51bca2b45daa47e289d69c09764}.
\end{center}
The catalogue at \texttt{docs/README.md}, the foundation report's
\texttt{README.md}, and its opening definitions were examined directly.
The catalogue maps seven surcomplex packages: analysis, analytic geometry,
finite deformations, contours and Stokes, global divisors, polynomial
algebra, and trigonometry. It explicitly separates coherent common-domain
sections from coefficientwise germs and warns against transferring theorems
between these rings~\cite{Repo}.

The topic selected here is holomorphic \emph{iteration and linearization},
not another treatment of local inversion, residues, divisor realization,
or finite zero schemes. In particular, inversion of a function and
conjugation of a dynamical system to its derivative are different problems.
Our support recursion solves the latter, and the exact radius-collapse
family shows that its common-domain answer can fail even though every
coefficient-germ step succeeds.

The new arguments do not depend on accepting every theorem asserted in
the repository's unrefereed drafts. The support, composition, homological,
and radius arguments needed here are given explicitly. The catalogue
inspection is not a claim to have performed a line-by-line audit of every
source manuscript or every Lean file in the repository.

\subsection{Classical local dynamics}

The ordinary homological denominators and recursive linearization methods
are classical. Marmi's lectures provide a systematic entry point, and
Yoccoz's work supplies the sharp quadratic obstruction used in the
illustrative example~\cite{Marmi,Yoccoz}. We do not reprove Yoccoz's theorem
or claim that its arithmetic threshold has been improved within its
original analytic category.

Carminati and Marmi study regular dependence of ordinary linearization on
the multiplier, including Whitney-smooth and monogenic
behavior~\cite{CarminatiMarmi}. Our reciprocal jets are related to this
natural parameter question, but our main statement concerns a different
requirement: after replacing two independent parameters by Hahn monomials,
do the resulting ordinary coefficient functions have a common disk? We
impose no Whitney estimates or joint analytic parameter convergence.

Fauvet, Menous, and Sauzin develop explicit tree expansions and recover
classical convergence estimates as well as refined multiplier
regularity~\cite{FauvetMenousSauzin}. The finite-dependence structure of our
recursion is compatible with such expansions; the use of recursive or
tree-organized formal coefficients is not itself a novelty claim.
Bounemoura's analysis of linear versus nonlinear arithmetic conditions
in different regularity classes gives useful conceptual context, but does
not concern the Hahn coefficient categories defined here~\cite{Bounemoura}.

\subsection{Non-Archimedean and Hahn analysis}

Lindahl studies analytic linearization disks over complete ultrametric
fields of equal characteristic zero, explicitly including Laurent series
fields with trivially valued complex coefficient fields~\cite{Lindahl}.
In that valuation, a nonzero complex small divisor is a unit of absolute
value one. Our quantity $\sigma$ instead measures its ordinary complex
absolute value, and therefore the analyticity of the \emph{coefficient
functions} in an ordinary variable. The common-domain obstruction is
invisible to the Hahn valuation alone. Proposition~\ref{prop:monad}
explains why unrestricted formal evaluation on the infinitesimal monad can
still work when ordinary coefficient analyticity fails.

M\"uller and Strohmaier develop Hahn-meromorphic functions with convergent
generalized power expansions on the logarithmic cover of the complex
plane~\cite{MuellerStrohmaier}. Their theory is an important precedent for
analytic theorems involving Hahn-type expansions. Its convergence
conditions and analytic variable differ from the formal positive-Hahn
parameter and ordinary coefficient-domain conditions used here. No theorem
is transferred between these categories without a separate proof.

\subsection{The novelty claim, delimited}

The specific candidate contributions are the exact multiplier-jet radius
law in Theorem~\ref{thm:collapse}, the resulting sharp nonlinear
common-domain/coefficient-germ trichotomy in
Theorem~\ref{thm:trichotomy}, and their multivariable counterpart in
Theorem~\ref{thm:multicollapse}. The positive lifting theorem supplies
the necessary arbitrary-support and actual-halo foundations. The sharp
obstruction uses only two positive monomials and rational input
coefficient functions.

The searched and consulted sources did not reveal these exact statements
in these Hahn coefficient categories. This is a provisional literature
assessment, not proof of priority. The linear homological estimates,
formal nonresonant conjugacy, support lemma, quadratic scaling trick,
and centralizer argument have classical antecedents and are not offered
as independent discoveries. A specialist could identify an existing
formal-family theorem that subsumes part of the present formulation;
that possibility should be checked before any publication claim.

\section{Reproducibility and exact finite verification}\label{sec:verification}

\subsection{What was computed}

The accompanying program \texttt{code/verify.py} uses only the Python
standard library and performs all arithmetic exactly in $\Q(i)$ using
\texttt{fractions.Fraction}. Its principal test multiplier is
\[
 \lambda=\frac35+\frac45i.
\]
This is not a root of unity: a root of unity in $\Q(i)$ is an algebraic
integer, hence lies in $\Z[i]$, whereas this number does not. This
argument is only used to justify nonzero denominators in exact test
examples; it is not an assertion about its continued-fraction loss exponent.

The recorded run used Python 3.13.5 and passed all 1,109 checks:
\begin{center}
\begin{tabular}{@{}lr@{}}
\toprule
Identity checked & Exact assertions\\
\midrule
Reciprocal-jet formula versus inverse recurrence &133\\
Truncated product of a denominator and its reciprocal &133\\
Ordinary quadratic Schr\"oder identity through degree 12 &13\\
Hahn-scaled quadratic identity &156\\
Nonlinear recursion with multiplier drift and cubic forcing &154\\
Two-variable multiplicative-detuning reciprocal jets &520\\
\midrule
Total &1109\\
\bottomrule
\end{tabular}
\end{center}
The first two groups use $2\le n\le20$ and $0\le k\le6$.
The mixed-forcing test uses
$F(z)=(\lambda+t)z+t(z^2+z^3)$ and checks the conjugacy after each
parameter step through $t^4$ and spatial degree ten. The multivariable
check uses the two exact unit phases
$(3+4i)/5$ and $(5+12i)/13$, total degrees two through nine, both components,
and reciprocal-jet orders zero through four. It checks nonzero
denominators in that finite range; it does not certify nonresonance for
all multiindices for this sample tuple.

For a readily inspectable output example, the quadratic coefficients are
\[
 b_2=\frac{11}{10}+\frac15 i,
 \qquad b_3=\frac{19}{20}+\frac{41}{40}i.
\]
The complete recorded output is in \texttt{data/verification.txt}.
To reproduce it, run
\begin{verbatim}
python code/verify.py
\end{verbatim}
from the article directory. The script raises an exception immediately if
an identity fails and prints the count of successful exact assertions.

\subsection{What the computation does not establish}

These finite checks are tests of algebraic formulas, signs, normalizations,
and the implementation of truncation. They do not prove Hahn summability,
well-founded recursion, ordinary analytic domain preservation, an infinite
small-divisor limsup, exact convergence radii, or novelty. Those conclusions
are supported by the mathematical arguments in this article, not by a
numerical experiment. No Lean theorem corresponding to the new analytic
results has been compiled or checked as part of this package.

\subsection{An implementable coefficient procedure}

For finitely generated positive support, a symbolic implementation can
follow \eqref{eq:recursion} directly. The input is a requested exponent
$\gamma$, the input functions indexed by $A$, and an oracle or certificate
for their ordinary homological inverse in the selected coefficient algebra.
First enumerate the finite words in $A$ summing to $\gamma$. Their letters
and partial sums give the finite dependency set. Then evaluate lower
coefficients in increasing exponent order, form the finite differential
right side, and divide each Taylor coefficient by its nonresonant divisor.
Finally, attach the inherited or computed domain certificate to the output.

This procedure is a mathematical finite-dependence description, not a
claim that arbitrary well-ordered supports or arbitrary holomorphic
functions have computable presentations. A computer algebra implementation
needs effective support and coefficient representations. In the polynomial
case, all operations are finite exact polynomial operations. In the
rational-function case, the homological inverse may cease to be rational,
so an exact implementation should retain a coefficient-generating operator
or a series object rather than pretending rational closure.

The obstruction family illustrates why a domain certificate cannot be a
single global Boolean flag. For its $u^kv$ coefficient, the appropriate
certificate is the exact radius $Re^{-(k+1)\sigma}$; reporting merely
``each coefficient converges'' would erase the decisive failure of a
common analytic neighborhood.

\section{Scope boundaries and further questions}\label{sec:scope}

\subsection{Fixed multiplier versus infinitesimal drift}

The sharp negative common-domain example varies the multiplier by $u$.
For the narrower class $F'(0)=\lambda$ exactly, the sufficiency of
$\sigma=0$ remains true by Theorem~\ref{thm:lifting}, and $\sigma=\infty$
can already obstruct coefficient germs at first order. The present argument
does not settle the optimal common-domain condition for all perturbations
with fixed multiplier when $0<\sigma<\infty$. Eliminating drift removes
the direct reciprocal-jet mechanism. A different counterexample, or a
positive cancellation theorem, would be needed for that narrower problem.

\subsection{Restricting the value group}

All positive results hold over every set-sized ordered abelian group.
The no-common-neighborhood counterexample is established for
$\Gstar=\Q+\sqrt2\Q$, and hence for the universal class of workspaces.
It also embeds in any surcomplex workspace containing this ordered subgroup.
The proof does not assert the same necessity statement for a fixed cyclic
value group such as $\Z$. There, different parameter blocks can collide
at the same Hahn exponent. Resolving the effect of those collisions is a
separate problem, not a consequence of the two-independent-exponent proof.

For the weaker requirement of retaining the \emph{original} disk, however,
one positive parameter is already enough: Proposition~\ref{prop:homological}
and the first-order equation show failure whenever $\sigma>0$.
The unresolved cyclic-group issue concerns the existence of \emph{some}
common smaller neighborhood for every input, not this weaker statement.

\subsection{Normed coefficient classes and ordinary parameter convergence}

The theorem imposes no common bound on $\norm{h_\gamma}_r$ as $\gamma$
varies. Adding weighted coefficient norms, Gevrey bounds in the parameter,
or ordinary convergence in $(u,v)$ produces stronger problems.
The compact estimate \eqref{eq:compactestimate} is a starting point for
such questions, but the present support recursion alone does not solve
them. Their arithmetic thresholds could differ from the common-domain
threshold established here.

\subsection{Resonance and global scales}

Resonant multipliers require normal forms with retained resonant terms.
Moreover, extending a finite-halo conjugacy to all infinite scales requires
compatibility conditions that ordinary coefficientwise entire functions do
not automatically satisfy. The critical-scale computation in
Section~\ref{sec:quadratic} supplies a concrete test that such an extension
must pass. No global conjugacy on all of $\No[i]$, no surreal-time flow,
and no arbitrary class-sized analytic construction is claimed.

\subsection{Conclusions}

The results answer a specific nontrivial question prompted by the
repository's distinct analytic rings: when does nonlinear infinitesimal
linearization stay in the common-domain ring, rather than merely existing
coefficient by coefficient? For irrational rotations and arbitrary
positive Hahn perturbations with multiplier drift, the answer is exact:
subexponential small divisors preserve the original domain; finite
exponential losses preserve individual germs but can destroy every common
domain; superexponential losses can destroy coefficient germs themselves.
The two-monomial reciprocal-jet construction measures that destruction
by the explicit law $Re^{-(k+1)\sigma}$, and the same mechanism persists
in several variables.

This is a category-sensitive extension of local dynamics, not a claim that
infinitesimals abolish the classical small-divisor problem. They separate
two issues that ordinary analytic convergence often forces one to solve
at once: algebraic dependence across perturbation orders and ordinary
analyticity within each coefficient. The sharp trichotomy describes the
resulting separation.

\appendix
\section{A dependency and hypothesis audit}\label{app:audit}

The new results can be checked along a short dependency chain.
The support lemma gives finite dependence at every Hahn exponent.
The homological inverse gives the analytic-space mapping property.
Their combination proves nonlinear lifting. The reciprocal-jet formula,
highest-pole dominance, and Cauchy--Hadamard then give an independent sharp
counterexample, which closes the universal characterization. The
multivariable proof repeats this chain with diagonal monomial divisors.

For a mathematical audit, the most consequential hypotheses are the
following. Positivity of the perturbation support makes all dependencies
strictly lower; arbitrary signed supports would invalidate that inference.
Well ordering ensures the sums exist and gives a least exponent in the
uniqueness argument. Vanishing at zero guarantees preservation of the
fixed point and the Taylor-order count. Nonresonance makes the homological
inverse exist. Unit-modulus reduction preserves the chosen disk or
polydisk. Common-domain analyticity is explicitly stronger than
coefficient-germ analyticity. Finally, rational independence of the two
obstruction exponents isolates the first nonlinear parameter block.

No compactness principle, limit of partial Hahn sums, or exchange of an
infinite ordinary sum with a Hahn coefficient extraction is hidden in the
proof. The limits used for radii are ordinary complex coefficient-growth
limits at a fixed Hahn exponent. The well-founded support recursion itself
is not a convergence process. This distinction is especially important
at \eqref{eq:critical}, where the naive Hahn sum is not defined at a
nonzero ordinary argument.

\section{A compact statement suitable for reuse}\label{app:reusable}

Let $B$ be one of the ordinary coefficient algebras used in this article,
closed under differentiation, finite products, and rotation by $\lambda$.
Assume the nonresonant homological inverse preserves $z^2B$.
Then every positive Hahn perturbation
$F=\lambda z+\sum t^a f_a$ with $f_a\in zB$ admits a unique normalized
linearizer with coefficients in $B$ and support in the additive monoid
generated by the input support. Its inverse has the same properties.
This assertion is exactly the proof of Theorem~\ref{thm:lifting} with
its coefficient algebra left abstract; it requires no norm on $B$.

The common-disk, germ, entire, and polynomial theorems are obtained by
checking this one closure property. The obstruction theorem shows that
such a closure check is sufficient but that passing to an inductive limit
of domains cannot simply be interchanged with taking a full Hahn family.
For finite positive $\sigma$, every coefficient belongs to the germ
inductive limit, while the entire Hahn family does not belong to the
common-domain inductive limit. This noncommutation is the algebraic form
of the exact radius-collapse phenomenon.

\begin{thebibliography}{99}

\bibitem{Repo}
Vladimir Reshetnikov, \emph{Surreal}, GitHub repository, inspected at
commit \texttt{39f2be6667ad} (full identifier in Section~\ref{sec:context}).
Catalogue \texttt{docs/README.md}; foundation material under
\texttt{docs/surcomplex/analysis/}.
\href{https://github.com/VladimirReshetnikov/Surreal/tree/39f2be6667ade51bca2b45daa47e289d69c09764}{Pinned repository snapshot}.
Accessed 21 September 2026. AI-assisted, unrefereed repository reports.

\bibitem{Higman}
Graham Higman, \emph{Ordering by divisibility in abstract algebras},
Proceedings of the London Mathematical Society (3) \textbf{2} (1952),
326--336. \doi{10.1112/plms/s3-2.1.326}.

\bibitem{Marmi}
Stefano Marmi, \emph{An Introduction to Small Divisors Problems},
Istituti Editoriali e Poligrafici Internazionali, 2000.
\arxiv{math/0009232}.

\bibitem{Yoccoz}
Jean-Christophe Yoccoz, \emph{Th\'eor\`eme de Siegel, nombres de Bruno et
polyn\^omes quadratiques}, in \emph{Petits diviseurs en dimension 1},
Ast\'erisque \textbf{231} (1995), 1--88 (Numdam pagination).
\href{https://www.numdam.org/item/AST_1995__231__1_0/}{Numdam publication record}.

\bibitem{CarminatiMarmi}
Stefano Marmi and Carlo Carminati,
\emph{Linearization of germs: regular dependence on the multiplier},
Bulletin de la Soci\'et\'e Math\'ematique de France \textbf{136} (2008),
no.~4, 533--564. \doi{10.24033/bsmf.2565}; \arxiv{0801.2844}.

\bibitem{FauvetMenousSauzin}
Fr\'ed\'eric Fauvet, Fr\'ed\'eric Menous, and David Sauzin,
\emph{Explicit linearization of one-dimensional germs through
tree-expansions}, Bulletin de la Soci\'et\'e Math\'ematique de France
\textbf{146} (2018), no.~2, 241--285.
\doi{10.24033/bsmf.2757}.

\bibitem{Bounemoura}
Abed Bounemoura, \emph{Optimal linearization of vector fields on the torus
in non-analytic Gevrey classes}, 2020. \arxiv{2006.05673}.

\bibitem{Lindahl}
Karl-Olof Lindahl, \emph{Linearization in ultrametric dynamics in fields
of characteristic zero---equal characteristic case}, p-Adic Numbers,
Ultrametric Analysis and Applications \textbf{1} (2009), no.~4, 307--316.
\arxiv{1111.1993}.

\bibitem{MuellerStrohmaier}
J\"orn M\"uller and Alexander Strohmaier, \emph{The theory of
Hahn-meromorphic functions, a holomorphic Fredholm theorem, and its
applications}, Analysis \& PDE \textbf{7} (2014), no.~3, 745--770.
\doi{10.2140/apde.2014.7.745}; \arxiv{1205.0236}.

\end{thebibliography}
\end{document}
